\section{Degree-test diagnostics of finite-network local heterogeneity}
\label{sec:degree-local-diagnostics}

The limiting fixed-count prediction has no static disorder at this
away-from-onset control point. A finite network may nevertheless have
distributed temporal variances and correlated changes of its local input
mean and variance. A positive input-mean spread must first be separated from
finite-record noise; it cannot be assigned to a new stationary mechanism by
looking only at a histogram of trajectory means.

For recording block \(k\), define the neuron-wise input mean and second
moment \(\mu_i^k,s_i^k\), rate mean and second moment \(m_i^k,t_i^k\), and
average gain \(g_i^k\). For two disjoint independent stationary blocks,
\begin{align*}
 \widehat q_i^{kl}
 &=\tfrac12(s_i^k+s_i^l)-\mu_i^k\mu_i^l,\\
 \widehat v_i^{kl}
 &=\tfrac12(t_i^k+t_i^l)-m_i^k m_i^l,\\
 \widehat A_i^{kl}&=g_i^kg_i^l.
\end{align*}
Their expectations are the infinite-time input variance, rate variance, and
squared time-averaged gain of the fixed network. The variance estimators
equal the diagonals of the mean-free two-window covariance construction.
No Gaussian assumption is required for these identities. Independence of
the windows is the substantive temporal assumption.

Let a tilde denote subtraction of a block's population mean over neuron
labels. Cross-block spatial products give
\[
 \widehat\Delta_a^{kl}
 =\frac1{N_a}\sum_{i\in a}\widetilde\mu_i^k\widetilde\mu_i^l,\qquad
 \widehat V_{m,a}^{kl}
 =\frac1{N_a}\sum_{i\in a}\widetilde m_i^k\widetilde m_i^l.
\]
For four distinct independent windows, products
\[
 \frac1{N_a}\sum_{i\in a}\widehat v_i^{kl}\widehat v_i^{mn},
 \quad
 \frac1{N_a}\sum_{i\in a}
  (\widehat q_i^{kl}-\overline{\widehat q}_a^{kl})
  (\widehat q_i^{mn}-\overline{\widehat q}_a^{mn})
\]
estimate \(E_i[v_i^2]\) and \(\operatorname{Var}_i q_i\), respectively.
A third independent window gives the local covariance
\[
 \widehat{\operatorname{Cov}}_a(\mu,q)
 =\frac1{N_a}\sum_{i\in a}
  \widetilde\mu_i^k
  (\widehat q_i^{lm}-\overline{\widehat q}_a^{lm}).
\]
The implementation averages all three matchings of each valid quadruple,
or all three choices of the singled-out window in each valid triple.
Gap changes test residual temporal dependence. None of these operations
produces another independent network replicate.

The estimated diagonal rate energy is
\(\widehat{\mathcal E}_{\rm diag}=\sum_a f_a\widehat E_i[v_i^2]\).
For nonzero denominators,
\[
 \widehat d=
 \left(\frac{\widehat\nu^2}{\widehat{\mathcal E}_{\rm diag}}\right)
 \left(\frac{\widehat{\mathcal E}_{\rm diag}}
             {\widehat{\mathcal E}}\right).
\]
This is an exact algebraic factorization of the estimates, not an assertion
that either ratio is unbiased or that a finite-record factor must lie in
its infinite-time range. Negative noise-corrected moments are retained.

In the prespecified retained \(N=1000,\chi=0\), seed4001 example, cross-block
mean products give
\[
 \widehat\Delta=(0.06067,0.09650),\qquad
 \widehat V_m=(0.003857,0.004693).
\]
The raw full-record mean variances are \((0.06368,0.10104)\).
Thus the temporal mean-estimation term is small compared with the observed
static spread. Substituting measured \(V_m\) and population means into the
exact finite-pool row contraction gives \((0.05939,0.09575)\), within about
two percent of the observed spread in this example. This substitution is
a calibrated diagnostic. Since the measured means depend on the same
network, the independent cavity-row contraction is not an exact identity
for that particular fixed connectivity.

The same cross-window calculation measures
\[
 \operatorname{Var}_i q_i=(0.2076,0.5473),\qquad
 \operatorname{Cov}_i(\mu_i,q_i)=(0.01962,0.05572).
\]
A local Gaussian mean-rate expansion identifies why these are relevant.
With
\[
 {\cal M}(\mu,q)=2\Phi\!\left(\frac{\mu}{\sqrt{1+q}}\right),
 \qquad
 \alpha=\partial_\mu{\cal M},\qquad
 \beta=\partial_\mu^2{\cal M},
\]
the Gaussian heat identity gives \(\partial_q{\cal M}=\beta/2\).
Expanding about a population background,
\[
 \delta m_i=\alpha\,\delta\mu_i+\frac{\beta}{2}\delta q_i+
             O(\delta\mu_i^2,\delta\mu_i\delta q_i,\delta q_i^2),
\]
and squaring yields
\[
 \operatorname{Var}_i m_i
 \simeq\alpha^2\Delta+
       \alpha\beta\operatorname{Cov}_i(\mu_i,q_i)+
       \frac{\beta^2}{4}\operatorname{Var}_i q_i.
\]
At the leading fixed-count background, the predicted and measured rate-mean
variances are
\begin{align*}
 V_m^{\rm Taylor}&=(0.004286,0.005276),\\
 \widehat V_m&=(0.003857,0.004693).
\end{align*}
The prediction exceeds the measured values by \(11\)--\(12\%\).
Using the measured population input mean and corrected temporal variance
reduces the excess to about \(7\%\). The terms involving variance
heterogeneity account for roughly a third of the expanded variance.

This diagnostic explains how finite-network temporal-variance heterogeneity
can feed static rate heterogeneity even when the leading fixed-count
solution is homogeneous within each population. It does not predict
\(\operatorname{Var}q\), \(\operatorname{Cov}(\mu,q)\), their scaling
coefficients, or the omitted nonlinear moments from microscopic parameters.
It is therefore not a controlled repair of the finite-size PR theory.
The \(N=2000\) checks and the full cohort results are recorded separately;
a selected example alone does not establish a size law.
